\section*{Capítulo \ref{ch:b1:intermolecular-forces-solvents} --- Forças intermoleculares e solventes}

\begin{solution}{exo:b1:intermolecular-forces-solvents:1}
Argônio: só London. Cloreto de hidrogênio: Keesom, Debye e London, com
predomínio de London (\cref{exo:b1:intermolecular-forces-solvents:8}); o
cloro é grande demais para \omterm{def:b1:intermolecular-forces-solvents:hbond}{ligações de hidrogênio} fortes. Água:
predominam as \omterm{def:b1:intermolecular-forces-solvents:hbond}{ligações de hidrogênio}, com os três termos de van der
Waals. Metano: só London. Di-iodo: só London, forte porque a molécula é
grande e muito polarizável.
\end{solution}

\begin{solution}{exo:b1:intermolecular-forces-solvents:2}
Entre suas próprias moléculas: \ce{CH3OH} (\ce{O-H}), \ce{CH3NH2}
(\ce{N-H}) e \ce{HF}. \ce{CH3OCH3} e \ce{CH3F} têm \omterm{def:b1:lewis-resonance-vsepr:lewis}{pares não ligantes} em
O ou F, mas nenhum hidrogênio ligado a eles: só podem aceitar \omterm{def:b1:intermolecular-forces-solvents:hbond}{ligações de
hidrogênio}, da água, por exemplo. \ce{CH4} não doa nem aceita.
\end{solution}

\begin{solution}{exo:b1:intermolecular-forces-solvents:3}
Água e etanol: polares, próticos. Propanona e acetonitrila: polares,
apróticas. Diclorometano: pouco polar, aprótico. Hexano: apolar, aprótico.
\end{solution}

\begin{solution}{exo:b1:intermolecular-forces-solvents:4}
O di-iodo é apolar: no hexano, ele troca \omterm{def:b1:intermolecular-forces-solvents:vdw}{interações de London} por
\omterm{def:b1:intermolecular-forces-solvents:vdw}{interações de London}, enquanto na água ele romperia \omterm{def:b1:intermolecular-forces-solvents:hbond}{ligações de
hidrogênio} sem substituí-las. O cloreto de sódio é iônico: só um
solvente de alta permissividade pode separar seus íons, e só um \omterm{def:b1:intermolecular-forces-solvents:solvent-classes}{solvente
polar} pode solvatá-los; o hexano ($\varepsilon_r = 1.89$) não faz nem uma
coisa nem outra.
\end{solution}

\begin{solution}{exo:b1:intermolecular-forces-solvents:5}
O metoximetano não tem \ce{O-H}: não pode doar \omterm{def:b1:intermolecular-forces-solvents:hbond}{ligações de hidrogênio} e
ferve à temperatura mais baixa. O metanol e o etanol formam ambos
\omterm{def:b1:intermolecular-forces-solvents:hbond}{ligações de hidrogênio}; o etanol, com um \ce{CH2} a mais, tem \omterm{def:b1:intermolecular-forces-solvents:vdw}{interações
de London} mais fortes. O ácido etanoico forma dímeros ligados por
\omterm{def:b1:intermolecular-forces-solvents:hbond}{ligações de hidrogênio} e é o mais pesado: ferve à temperatura mais alta.
\end{solution}

\begin{solution}{exo:b1:intermolecular-forces-solvents:6}
No vácuo, $F = e^2/(4\pi\varepsilon_0 r^2) = (1.602 \times
10^{-19})^2/(4\pi \times 8.854 \times 10^{-12} \times (5.00 \times
10^{-10})^2) = \qty{9.23e-10}{N}$. Na água, $F/78.4 = \qty{1.18e-11}{N}$;
no hexano, $F/1.89 = \qty{4.88e-10}{N}$. A água enfraquece a atração 41
vezes mais que o hexano: ela consegue manter os íons separados, o hexano não.
\end{solution}

\begin{solution}{exo:b1:intermolecular-forces-solvents:7}
Misturar propanona com água rompe algumas \omterm{def:b1:intermolecular-forces-solvents:hbond}{ligações de hidrogênio}
água--água, mas forma novas ligações entre a água e o oxigênio de
\ce{C=O}, além de interações dipolo--dipolo: o balanço é favorável. O
hexano só oferece à água fracas \omterm{def:b1:intermolecular-forces-solvents:vdw}{interações de London} em troca das
\omterm{def:b1:intermolecular-forces-solvents:hbond}{ligações de hidrogênio} que ela romperia: os dois líquidos ficam separados.
\end{solution}

\begin{solution}{exo:b1:intermolecular-forces-solvents:8}
\ce{HBr} tem mais elétrons e uma nuvem maior e mais polarizável: sua
\omterm{def:b1:intermolecular-forces-solvents:vdw}{interação de London} supera a de \ce{HCl} mais do que sua \omterm{def:b1:intermolecular-forces-solvents:vdw}{interação de
Keesom} fica abaixo dela. London predomina.
\end{solution}

\begin{solution}{exo:b1:intermolecular-forces-solvents:9}
\chemfig{CH_3-C(=[:60]O)-[:-60]O-[:0]H} frente a seu parceiro, com duas
ligações $\ce{O-H}\cdots\ce{O=C}$ fechando um anel de oito membros. No
benzeno, apolar e aprótico, nada compete com essas ligações: o ácido está
dimerizado e se comporta como uma partícula de cerca de \qty{120}{g/mol}.
Na água, cada molécula de ácido forma \omterm{def:b1:intermolecular-forces-solvents:hbond}{ligações de hidrogênio} com a água:
os dímeros se desfazem.
\end{solution}

\begin{solution}{exo:b1:intermolecular-forces-solvents:10}
O cloreto de sódio é formado de íons: a água, de alta permissividade,
os separa (dissociação) e os hidrata (oxigênio voltado para \ce{Na+},
hidrogênio voltado para \ce{Cl-}). O cloreto de hidrogênio é uma \omterm{def:b1:lewis-resonance-vsepr:dipole}{molécula
polar}: a água, polar, transforma a ligação \ce{H-Cl} em um par de íons
\ce{H3O+ Cl-} (ionização), separa o par (dissociação) e hidrata os íons.
O hexano é apolar e de baixa permissividade: não ioniza nem dissocia, de
modo que não há íons para conduzir a corrente.
\end{solution}

\begin{solution}{exo:b1:intermolecular-forces-solvents:11}
$(174.1 - 36.1)/5 = \qty{27.6}{\celsius}$ por \ce{CH2}. O último
incremento (do nonano ao decano) é de \qty{23.4}{\celsius}, de modo que o
undecano deveria ferver perto de $174 + 22 \approx \qty{196}{\celsius}$.
Cada \ce{CH2} acrescenta a mesma contribuição de London, mas uma fração
cada vez menor do total da molécula, enquanto a temperatura necessária
para vencer as interações já é alta.
% ledger: bp:C9H20
\end{solution}

\begin{solution}{exo:b1:intermolecular-forces-solvents:12}
Parte \omterm{def:b1:intermolecular-forces-solvents:amphiphile}{hidrofílica}: a cabeça sulfato \ce{-OSO3^-} com seu contraíon \ce{Na+};
parte \omterm{def:b1:intermolecular-forces-solvents:amphiphile}{hidrofóbica}: a cadeia de doze carbonos. Na superfície, cabeças na
água e caudas no ar; em uma micela, caudas dentro e cabeças fora. As
caudas se dissolvem na gordura da mancha, as cabeças a mantêm em contato
com a água, e a gordura é removida em micelas.
\end{solution}

\begin{solution}{pb:b1:intermolecular-forces-solvents:1}
\textbf{1.} A \omterm{def:b1:intermolecular-forces-solvents:vdw}{interação de London}, a única entre átomos apolares.
\textbf{2.} Do hélio ao xenônio, o número de elétrons e o tamanho
crescem, e com eles a \omterm{def:b1:periodicity:polarisability}{polarizabilidade} e a atração de London.
\textbf{3.} O hélio tem dois elétrons fortemente retidos pelo núcleo: é o
átomo menos polarizável, com a atração de London mais fraca.
\textbf{4.} $68.8 - 36.1 = \qty{32.7}{\celsius}$: um \ce{CH2} acrescenta
elétrons e superfície de contato, portanto atração de London.
\textbf{5.} London, que cresce com a \omterm{def:b1:periodicity:polarisability}{polarizabilidade}: $\ce{CH4} < \ce{SiH4} <
\ce{GeH4}$.
\textbf{6.} O carbono não é eletronegativo o bastante para polarizar
fortemente \ce{C-H}, e o metano não tem \omterm{def:b1:lewis-resonance-vsepr:lewis}{par não ligante} para aceitar uma
\omterm{def:b1:intermolecular-forces-solvents:hbond}{ligação de hidrogênio}.
\textbf{7.} A água: polar, prótica, com a maior permissividade; ela
separa e solvata os dois íons.
\textbf{8.} $78.4/1.89 = 41$: a atração é 41 vezes mais fraca na água.
\textbf{9.} A acetonitrila (ou a propanona): polar o bastante para
dissolver o sal, aprótica para que o ânion fique livre.
\textbf{10.} O hexano: imiscível com a água, apolar como o di-iodo. O iodo
passa para o hexano, a fase superior (o hexano é menos denso que a água).
\textbf{11.} O etanol é \omterm{def:b1:intermolecular-forces-solvents:miscibility}{miscível} com a água: não se forma uma segunda fase.
\textbf{12.} Seu oxigênio aceita \omterm{def:b1:intermolecular-forces-solvents:hbond}{ligações de hidrogênio} da água, mas ele
não pode doar nenhuma, e seus dois grupos etila são \omterm{def:b1:intermolecular-forces-solvents:amphiphile}{hidrofóbicos}: ele só
se dissolve um pouco.
\textbf{13.} $-66.4 - (-85.0) = \qty{18.6}{\celsius}$ por \omterm{def:b1:periodicity:table}{período};
\ce{HF} extrapolado: $-85.0 - 18.6 = \qty{-103.6}{\celsius}$.
\textbf{14.} \ce{HF} ferve a \qty{19.6}{\celsius}, \qty{123}{\celsius}
acima da extrapolação.
\textbf{15.} $-62.5 - (-87.8) = \qty{25.3}{\celsius}$; \ce{NH3}
extrapolado: \qty{-113.1}{\celsius}; valor real $-33.4$, \qty{80}{\celsius}
mais alto.
\textbf{16.} $-88.1 - (-112.0) = \qty{23.9}{\celsius}$; \ce{CH4}
extrapolado: \qty{-135.9}{\celsius}; valor real $-161.5$, abaixo. O metano
não forma \omterm{def:b1:intermolecular-forces-solvents:hbond}{ligações de hidrogênio}; ele é simplesmente pequeno e pouco
polarizável.
\textbf{17.} Água (cerca de \qty{179}{\celsius}) $>$ \ce{HF} (123) $>$
\ce{NH3} (80).
\textbf{18.} \ce{NH3} pode doar três, mas só aceitar uma (um \omterm{def:b1:lewis-resonance-vsepr:lewis}{par não
ligante}); \ce{HF} pode aceitar três, mas só doar uma. A água doa duas e
aceita duas: cada molécula pode participar de quatro \omterm{def:b1:intermolecular-forces-solvents:hbond}{ligações de
hidrogênio}, uma rede tridimensional completa.
\textbf{19.} \ce{H2S} e \ce{H2Se} não formam \omterm{def:b1:intermolecular-forces-solvents:hbond}{ligações de hidrogênio}
significativas: suas temperaturas de ebulição seguem a tendência de
London que a água seguiria sem elas. A própria água é a anomalia que se
quer medir.
\textbf{20.} $-41.3 - (-60.3) = \qty{19.0}{\celsius}$ por \omterm{def:b1:periodicity:table}{período}.
\textbf{21.} $T(p) = -60.3 + 19.0\,(p - 3)$ em \unit{\celsius}.
\textbf{22.} $T(2) = -60.3 - 19.0 = \qty{-79.3}{\celsius}$.
\textbf{23.} $100.0 - (-79.3) = \qty{179}{\celsius}$.
\textbf{24.} Sem suas \omterm{def:b1:intermolecular-forces-solvents:hbond}{ligações de hidrogênio}, a água ferveria perto de
\textbf{\qty{-79}{\celsius}}, e seria um gás na Terra.
\end{solution}
